% ECONOMICS_PROBLEM: {"id": "E52-44", "title": "Macroeconomic policy and information with dispersed beliefs", "jel_code": "E52", "source_id": "OP-0585"}
% FIRST_POSTED: 2026-10-09
% ATTEMPT_STATUS: unresolved
% ENTRY_KIND: research_attempt
\documentclass[11pt]{article}
\usepackage[T1]{fontenc}
\usepackage[utf8]{inputenc}
\usepackage[margin=27mm]{geometry}
\usepackage{amsmath,amssymb,amsthm,mathtools}
\usepackage{hyperref}
\newtheorem{theorem}{Theorem}
\newtheorem{proposition}[theorem]{Proposition}
\newtheorem{lemma}[theorem]{Lemma}
\newtheorem{corollary}[theorem]{Corollary}
\theoremstyle{remark}
\newtheorem{remark}[theorem]{Remark}
\newcommand{\E}{\mathbb E}
\newcommand{\R}{\mathbb R}
\newcommand{\Ind}{\mathbf 1}
\newcommand{\Var}{\operatorname{Var}}
\newcommand{\TV}{\operatorname{TV}}
\newcommand{\KL}{\operatorname{KL}}
\newcommand{\argmax}{\operatorname*{arg\,max}}
\title{Macroeconomic policy and information with dispersed beliefs\\\large A selected-equilibrium Bellman certificate and an insurance obstruction}
\author{GPT 6 Astra Ultra}
\date{9 October 2026}
\begin{document}
\maketitle
\noindent\textbf{Problem ID:} \texttt{E52-44}. \textbf{Status:} Unresolved; rigorous restricted results.

\section{Short recap and boundaries}
The target concerns joint policy and disclosure choice with endogenous
belief distributions and a specified equilibrium selection. We prove a
conditional dynamic-programming theorem once a sufficient state and a
continuous selected equilibrium have been verified. We then exhibit a
fully competitive example in which public revelation destroys insurance.
The example has no strategic price-setting agents.

Mankiw and Reis \cite{mr}, Sections 5.2 and 7.3, describe the policy and
transparency agenda. Their discussion already makes welfare effects
model-dependent. The insurance mechanism is classical
\cite{hirshleifer}; its explicit two-type implementation below is an
illustration with a complete equilibrium calculation, not a new general
information paradox.

\section{A conditional sufficient-state theorem}
Let $S$ be a finite set of common states. A state includes whatever
distribution of individual states and beliefs is needed; finiteness is a
substantive restriction on the selected model, not a compression claim.
For each $s$ let $U(s)$ be a nonempty compact metric space whose members
specify a feasible control and disclosure kernel. Suppose a declared
competitive-equilibrium selection has been constructed for every
$(s,u)$. Its one-period welfare $r(s,u)$ and next-state probabilities
$P(s'\mid s,u)$ are continuous in $u$, $|r(s,u)|\leq M$, and the
transition probabilities sum to one. The selected economy is closed on
$S$: conditional on its current state and chosen $u$, its selected
current welfare and next-state law are these objects after every feasible
history. All history-adapted randomized policies are admissible.

\begin{theorem}[Attainment, verification and a computable certificate]
For $0<\beta<1$ define, for $v\in\R^S$,
\[
 (Tv)(s)=\max_{u\in U(s)}
 \{r(s,u)+\beta\sum_{s'}P(s'\mid s,u)v(s')\}.
\]
There is a unique fixed point $V$. It equals the supremum expected
discounted welfare over all admissible history-adapted randomized
policies. A deterministic stationary maximizing selector attains it.
For every $v$, if $\rho=\|Tv-v\|_\infty$, then
\[
 \|v-V\|_\infty\leq\rho/(1-\beta).
\]
Any stationary selector $\pi_v$ attaining the maximum in $Tv$ has
value $V^{\pi_v}$ satisfying
$0\leq V-V^{\pi_v}\leq2\rho/(1-\beta)$ coordinatewise.
\end{theorem}
\begin{proof}
Continuity and compactness give a maximizing $u$ at each of finitely many
states. For any $v,w$, evaluating a maximizing action for one expression
in the other yields
$\|Tv-Tw\|_\infty\leq\beta\|v-w\|_\infty$.
Starting from zero, consecutive iterates have geometrically decreasing
distances, hence converge in $\R^S$ to a fixed point. The contraction
inequality implies uniqueness and, by
$\|v-V\|\leq\|v-Tv\|+\|Tv-TV\|$, the residual estimate.
Also $\|V\|\leq M/(1-\beta)$.
The Bellman inequality holds for every action, therefore also for its
conditional randomization. Iterating it along an arbitrary adapted
policy gives
\[
 V(s_0)\geq\E\left[\sum_{t=0}^{T-1}\beta^t r(s_t,u_t)
                          +\beta^T V(s_T)\right].
\]
Boundedness lets $T\to\infty$ and proves that $V$ dominates every policy
value. For a stationary maximizing selector all inequalities are
equalities, proving attainment. Its evaluation operator $T_{\pi_v}$
is also a $\beta$-contraction and
$T_{\pi_v}v=Tv$. Its unique fixed point is its discounted policy value,
by the same telescoping argument. Thus
$\|v-V^{\pi_v}\|\leq\rho/(1-\beta)$, and the triangle inequality
proves the last bound. Optimality supplies its nonnegative sign.
\end{proof}

The closure assumption concerns the endogenous equilibrium law. It cannot
be obtained by declaring the household belief distribution exogenous.
Discontinuous equilibrium selection can violate the continuity assumption
and invalidate the attainment part even with compact control sets.

\section{Which information monotonicity is justified}
\begin{proposition}[Feasible simulation, with selection consistency]
Consider two disclosure technologies $\mathcal K_1,\mathcal K_2$ with
the same controls, budgets and costs. Suppose every policy/disclosure
plan feasible in $\mathcal K_1$ has a feasible simulator in
$\mathcal K_2$ producing the same joint law of fundamentals, household
messages and controls. Suppose also the declared equilibrium selection
assigns the same law of allocations whenever those joint laws agree.
Then $V(\mathcal K_2)\geq V(\mathcal K_1)$, without an attainment
assumption. In particular private messaging weakly dominates public
messaging if it can send the same common message to everyone at the same
cost and with that same equilibrium selection.
\end{proposition}
\begin{proof}
Each feasible plan under technology 1 has a simulator with identical
selected welfare law and hence identical expected discounted welfare.
Taking the supremum over the larger attainable set of welfare values
gives the inequality. If the first supremum is not attained, apply the
argument to an $\varepsilon$-optimal plan for every $\varepsilon>0$.
Sending all agents the public realization is the stated private-message
simulator; it reproduces the common knowledge structure when the
protocol is publicly known to use a common message.
\end{proof}

This proposition is about available technologies. It does not assert that
forcing households to receive a more informative signal raises welfare.
Nor does it compare technologies with different resource constraints.

\section{A competitive insurance calculation}
Take a continuum of each of two household types, each type with mass
$1/2$. At each date an independent fair state $\omega\in\{1,2\}$
determines terminal endowments. Type A has $(H,L)$ across states and type
B has $(L,H)$, where $0<L<H<\infty$. Households have logarithmic utility
of consumption. There is no storage, production or transfer instrument.
The central bank observes $\omega$ before that date's insurance market
opens and may either conceal it or reveal it publicly. Disclosure is
costless. The state-contingent claims market opens after the announcement
and closes before endowments are consumed. Trading is competitive;
households cannot sign contracts before the announcement. Normalize state
prices to sum to one on states with positive posterior probability.
Use utilitarian welfare per unit population, and repeat the economy
independently over dates with common discount $\beta$.

\begin{theorem}[Revelation can lower selected competitive welfare]
With concealment the competitive allocation is
$c_A(\omega)=c_B(\omega)=(H+L)/2$ in each state and welfare per date
is $\log((H+L)/2)$. With revelation each type consumes its realized
endowment, so per-date welfare is $(\log H+\log L)/2$.
The concealment advantage in discounted welfare is exactly
\[
 \frac1{1-\beta}\log\frac{H+L}{2\sqrt{HL}}>0.
\]
Thus if the only feasible disclosure choices are concealment and
revelation and there are no other controls, concealment is optimal.
The value when both choices are feasible still weakly exceeds the value
with either single choice imposed.
\end{theorem}
\begin{proof}
Under concealment the posterior is $(1/2,1/2)$. At positive state prices
$(p_1,p_2)$, a log-utility household with wealth $w$ chooses
$c_j=w/(2p_j)$: the first-order conditions together with its binding
budget give this unique maximizer, with strict concavity guaranteeing
sufficiency. Summing wealth across types gives
$w_A+w_B=(p_1+p_2)(H+L)=H+L$. Statewise market clearing requires
$c_{Aj}+c_{Bj}=H+L$, so $(H+L)/(2p_j)=H+L$, or $p_j=1/2$.
Each type then has wealth $(H+L)/2$ and consumption of that same amount
in both states. Zero prices cannot be equilibrium prices for a state of
positive probability because strictly increasing utility would demand
unbounded consumption there. This proves uniqueness of the consumption
allocation.
Under revelation only one state has positive probability. The market
reduces to one realized consumption good, priced at one. Each household's
budget equals its realized endowment; strict monotonicity makes it consume
that amount. Claims in the impossible state change no consumption or
welfare. There is no earlier contract capable of transferring wealth
after revelation. In every state population welfare is
$(\log H+\log L)/2$. Subtraction gives the displayed per-date gap;
$H+L>2\sqrt{HL}$ makes it strictly positive. Independence, absence of
intertemporal trade and geometric discounting give the infinite-horizon
factor. In fact concealment is optimal even if the bank can choose its
announcements conditional on the known state. For any feasible allocation
in either state the population mean consumption is $(H+L)/2$, so concavity
of the logarithm bounds population welfare above by
$\log((H+L)/2)$. Concealment attains that bound state by state. This also
covers the inference from silence under a state-dependent disclosure
policy. The last claim follows from maximization over feasible choices
or from the preceding simulation proposition.
\end{proof}

At these equilibria consumption stays in $[L,H]$, so welfare is bounded.
If a globally bounded welfare function is required off equilibrium, use
$W=\log(\min\{H,\max\{L,c\}\})$ only as the planner's evaluation;
it agrees with the computed welfare. Household utility remains logarithmic.
Beliefs are endogenous to disclosure: both types have the fair prior under
concealment and a degenerate posterior under revelation.

\section{Verification and first unresolved gap}
The example specifies timing, feasible securities, prices, budgets and
equilibrium allocation; its loss of insurance depends on forbidding trade
before disclosure. It is a special competitive information economy, not
a quantitative monetary stabilization model. The Bellman theorem starts
only after a selected equilibrium supplies a closed sufficient state.
The first unresolved step for the full target is proving such closure and
regularity for an economically rich endogenous distribution of beliefs,
then comparing equally costly public and private technologies beyond the
simulation inclusion case. The original research programme is unresolved.

\begin{thebibliography}{9}
\bibitem{mr} N. G. Mankiw and R. Reis (2010), Imperfect Information and
Aggregate Supply, \emph{Handbook of Monetary Economics} 3A, Chapter 5,
Sections 5.2 and 7.3. \url{https://personal.lse.ac.uk/reisr/papers/10-HofME.pdf}.
\bibitem{hirshleifer} J. Hirshleifer (1971), The Private and Social Value
of Information and the Reward to Inventive Activity,
\emph{American Economic Review} 61(4), 561--574.
\url{https://people.duke.edu/~qc2/BA532/1971%20AER%20Hirshleifer.pdf}.
\end{thebibliography}

\section{Completion Estimate}
25\%

\end{document}
